\chapter{Dataset}\label{chapter:Dataset}

\section{Dataset Construction}
The base structure of the dataset consists of three columns: the website's HTML, 
a corresponding image, such as a screenshot, and the accessibility issues found 
in the HTML. \newline

\subsection{Data Collection}
To make sure that the following steps deliver a meaningful result, the 
dataset created for this study, tries to reflect real-world UI examples
from different areas. \newline
There have been a few attempts and analysis in the area of design-to-code, in 
order to see how well current LLMs can perform this task. Some Researches  
have open-sourced their dataset which allowed us to selectively pick good 
examples for this study. Especially, the dataset \textit{Design2Code}~\parencite{si2024design2code} 
and \textit{WebCode2M}~\parencite{gui2024webcode2m} have described promising 
approaches concerning the data collection and cleaning. Both datasets contain
real-world examples of webpages. \newline
Therefore, 28 examples of the dataset \textit{Design2Code} and 25 examples of 
the dataset \textit{WebCode2M} have been manually selected in order to ensure 
a wide range of different areas and styles. \newline
Manually erzählen

\subsection{Data Refinement}
Due to certain characteristics

\subsubsection{Code Purification}
Webpage code often contains redundant or not visible elements which can 
cause violations against certain wcag criteria. Since those elements are not 
required for image-to-code tasks, we have filtered them out the data entries
to make sure that our analysis provides a fair comparison of real-world and 
generated frontend code.

\subsubsection{No External Dependencies}
For our study, we assume a setting where only screenshots and corresponding 
prompts will be provided. This means that no external sources for elements 
such as images are taken into account. Instead of an external source, we 
only provide a \textit{placeholder} source without changing the general 
layout of the UI. \newline
For similar reasons, the datasets, stated above, have deleted all href 
attributes and links in the frontend code. While this does not change the 
layout of a webpage, it automatically violates against WCAG CRITERIA (GENAU).
To prevent the assumption that every webpage violates against this criteria, 
we have added empty href links to the dataset.


\subsection{Accessibility Violations}
Due to certain characteristics